\documentclass[10pt,a4paper]{amsart}
\usepackage[margin=19mm]{geometry}
\usepackage{amsmath,amssymb,mathtools}
\usepackage[scr=rsfs,cal=euler]{mathalfa}
\usepackage{xfrac}
\usepackage[unicode,hidelinks,bookmarks=false]{hyperref}
\newcommand{\ie}{i.e.\ }
\newcommand{\cf}{cf.\ }
\newcommand{\resp}{respectively\ }
\newcommand{\Subsection}{\subsection}
\providecommand{\nobreakdash}{\nobreak\hbox{-}}
\setlength{\emergencystretch}{2em}
\multlinegap0pt
\usepackage{amssymb}
\usepackage{mathtools}
\usepackage[shortlabels]{enumitem}
\newtheorem{lemma}{Lemma}
\newcommand{\dd}{\,\mathrm d}
\title{High-energy asymptotics for finite-interval Schr\"odinger operators with Gaussian white-noise potential}
\author{Yingdu Dong, Wenwen Jian, Xiaoping Yuan}
\date{Source: arXiv:2606.22426v1 (CC BY 4.0)}
\begin{document}
\maketitle

\noindent\emph{Source and licence.} High-energy asymptotics for finite-interval Schr\"odinger operators with Gaussian white-noise potential. Version arXiv:2606.22426v1 (CC BY 4.0), distributed by arXiv under the Creative Commons Attribution 4.0 International licence, http://creativecommons.org/licenses/by/4.0/, as recorded in the arXiv licence metadata for this exact version and shown on the abstract page https://arxiv.org/abs/2606.22426v1. Full licence notice: \texttt{licenses/arxiv-2606.22426-CC-BY-4.0.txt}.\par\smallskip

\noindent\textit{Editorial scope.} Counted unit: the article's lemma lem:phase-uniform (source0.tex lines 722--734). The article's complete original proof of it is delivered with it (source0.tex lines 736--925, one \texttt{proof} environment of 190 lines). No mathematics is written, repaired or re-derived: the statement and the proof are the article's own lines, in the article's own notation, and no retained line is substituted or re-flowed.  Retained uncounted as the source-local context the proof needs: lines 497--501; lines 609--612.  Nothing was omitted from any retained span, so the package carries no omission record.  No retained line contains a citation, so the wrapper carries no bibliography and no bibliography entry.  No retained line contains a figure environment, an image-inclusion command or drawing code, so no image is delivered and the package declares no asset dependency.  Editorial formatting only, in addition: an AMS \texttt{amsart} preamble that restates the article's own declarations and macros used by the retained lines, namely the environments \texttt{lemma} on separate counters, together with the standard packages those lines need.  The retained environments print in this document as Lemma 1 in order of appearance, whereas the article prints the same items with the numbering of its own full document.  The complete native source of the article as distributed in the arXiv e-print for this exact version is bundled unchanged as \texttt{sources/203346/source0.tex}.
\begin{equation}\label{eq:prufer-angle}
    \dd\theta_x
    =k\,\dd x-\frac{\rho}{k}\sin^2\theta_x\,\dd B_x
    +\frac{\rho^2}{k^2}\sin^3\theta_x\cos\theta_x\,\dd x.
\end{equation}

\begin{equation}\label{eq:delta-int}
    \delta_x=-\frac{\rho}{k}\int_0^xh(\theta_s)\,\dd B_s
    +\frac{\rho^2}{k^2}\int_0^xg(\theta_s)\,\dd s.
\end{equation}

\begin{lemma}[Parameter-uniform phase estimates]\label{lem:phase-uniform}
Let
\[
    I_n=[K_n-1,K_n+1],\qquad K_n=\frac{n\pi}{L}.
\]
For every finite $p\ge2$ there is a constant $C_p$, independent of $n$, such that, for all sufficiently large $n$,
\begin{align*}
    \left\|\sup_{k\in I_n}\sup_{x\le L}|\theta_x(k)-kx|\right\|_{L^p(\Omega)}
    &\le C_p n^{-1},\\
    \left\|\sup_{k\in I_n}\sup_{x\le L}|\partial_k\theta_x(k)-x|\right\|_{L^p(\Omega)}
    &\le C_p n^{-1}.
\end{align*}
\end{lemma}

\begin{proof}
We write the proof in some detail, since this is the point where one has to distinguish fixed-frequency BDG estimates from estimates which are uniform in the parameter $k$.

First fix a deterministic $k\ge1$ and put
\[
    \delta_x(k):=\theta_x(k)-kx .
\]
The integrated equation \eqref{eq:delta-int} gives
\[
    \delta_x(k)
    =-\frac{\rho}{k}\int_0^x h(\theta_s(k))\,\dd B_s
      +\frac{\rho^2}{k^2}\int_0^x g(\theta_s(k))\,\dd s .
\]
Since $h$ and $g$ are bounded, BDG implies
\[
\begin{aligned}
    \left\|\sup_{x\le L}|\delta_x(k)|\right\|_{L^p(\Omega)}
    &\le \frac{C_p}{k}
       \left\|\left(\int_0^L |h(\theta_s(k))|^2\,\dd s\right)^{1/2}\right\|_{L^p(\Omega)}
       +\frac{C}{k^2}\int_0^L\dd s  \\
    &\le C_p k^{-1}.
\end{aligned}
\]
This gives the first fixed-frequency bound.

Next set
\[
    \zeta_x(k):=\partial_k\theta_x(k),\qquad e_x(k):=\zeta_x(k)-x .
\]
The differentiability with respect to $k$ follows from the usual smooth parameter-dependence theorem for strong solutions of SDEs. In the estimates below we work on the compact intervals $I_n=[K_n-1,K_n+1]$; for all large $n$, these intervals lie in a fixed half-line $[k_0,\infty)$ with $k_0>0$. On such sets the coefficients
\[
    b(k,u)=k+\frac{\rho^2}{k^2}g(u),
    \qquad
    a(k,u)=-\frac{\rho}{k}h(u)
\]
are smooth in $(k,u)$, the relevant $k$-derivatives are locally bounded, and all derivatives in $u$ are bounded uniformly. Hence, after choosing a modification if necessary, $(k,x)\mapsto\theta_x(k)$ is twice continuously differentiable in $k$ on compact $k$-intervals, uniformly for $x\in[0,L]$ almost surely, and the $k$-derivatives solve the equations obtained by differentiating the stochastic equation. Differentiating \eqref{eq:prufer-angle} gives
\[
\begin{aligned}
    \dd\zeta_x
    &=\dd x+\frac{\rho}{k^2}h(\theta_x)\,\dd B_x
      -\frac{\rho}{k}h'(\theta_x)\zeta_x\,\dd B_x \\
    &\quad -\frac{2\rho^2}{k^3}g(\theta_x)\,\dd x
      +\frac{\rho^2}{k^2}g'(\theta_x)\zeta_x\,\dd x,
    \qquad \zeta_0=0.
\end{aligned}
\]
Subtracting $x$ and using $\zeta_s=s+e_s$, we obtain
\[
\begin{aligned}
    e_x
    ={}&\frac{\rho}{k^2}\int_0^x h(\theta_s)\,\dd B_s
      -\frac{\rho}{k}\int_0^x h'(\theta_s)s\,\dd B_s
      -\frac{\rho}{k}\int_0^x h'(\theta_s)e_s\,\dd B_s \\
    &-\frac{2\rho^2}{k^3}\int_0^x g(\theta_s)\,\dd s
      +\frac{\rho^2}{k^2}\int_0^x g'(\theta_s)s\,\dd s
      +\frac{\rho^2}{k^2}\int_0^x g'(\theta_s)e_s\,\dd s .
\end{aligned}
\]
Let
\[
    A_p(k):=\left\|\sup_{x\le L}|e_x(k)|\right\|_{L^p(\Omega)} .
\]
For the deterministic drift terms, boundedness of $g'$ gives
\[
    \left\|\sup_{x\le L}\left|\frac{\rho^2}{k^2}\int_0^x g'(\theta_s)s\,\dd s\right|\right\|_{L^p(\Omega)}
    \le Ck^{-2},
\]
Similarly, the term with $g$ is $O(k^{-3})$. For the last drift term,
\[
\begin{aligned}
    \left\|\sup_{x\le L}\left|\frac{\rho^2}{k^2}\int_0^x g'(\theta_s)e_s\,\dd s\right|\right\|_{L^p(\Omega)}
    &\le Ck^{-2}\int_0^L \|e_s\|_{L^p(\Omega)}\,\dd s  \\
    &\le Ck^{-2}A_p(k).
\end{aligned}
\]
For the stochastic terms, BDG gives
\[
    \left\|\sup_{x\le L}\left|\frac{\rho}{k^2}\int_0^x h(\theta_s)\,\dd B_s\right|\right\|_{L^p(\Omega)}
    \le C_p k^{-2},
\]
\[
    \left\|\sup_{x\le L}\left|\frac{\rho}{k}\int_0^x h'(\theta_s)s\,\dd B_s\right|\right\|_{L^p(\Omega)}
    \le \frac{C_p}{k}\left(\int_0^L s^2\,\dd s\right)^{1/2}
    \le C_p k^{-1},
\]
and
\[
\begin{aligned}
    \left\|\sup_{x\le L}\left|\frac{\rho}{k}\int_0^x h'(\theta_s)e_s\,\dd B_s\right|\right\|_{L^p(\Omega)}
    &\le \frac{C_p}{k}
      \left\|\left(\int_0^L |e_s|^2\,\dd s\right)^{1/2}\right\|_{L^p(\Omega)}  \\
    &\le C_p k^{-1}A_p(k).
\end{aligned}
\]
Combining these estimates yields
\[
    A_p(k)\le C_p k^{-1}+C_p(k^{-1}+k^{-2})A_p(k).
\]
Choose $k_0$ so large that $C_p(k^{-1}+k^{-2})\le1/2$ for $k\ge k_0$. Then the last term can be absorbed into the left-hand side, and hence
\begin{equation}\label{eq:e-fixed-bound}
    \left\|\sup_{x\le L}|\partial_k\theta_x(k)-x|\right\|_{L^p(\Omega)}
    =A_p(k)
    \le C_p k^{-1}.
\end{equation}

We shall also need the same size bound for the $k$-derivative of $e$. Put
\[
    \eta_x(k):=\partial_k^2\theta_x(k)=\partial_k e_x(k).
\]
Differentiating the equation for $\zeta$ once more gives
\[
\begin{aligned}
\dd\eta_x={}&
\left[-\frac{2\rho}{k^3}h(\theta_x)
      +\frac{2\rho}{k^2}h'(\theta_x)\zeta_x
      -\frac{\rho}{k}h''(\theta_x)\zeta_x^2
      -\frac{\rho}{k}h'(\theta_x)\eta_x\right]\dd B_x \\
&+\left[\frac{6\rho^2}{k^4}g(\theta_x)
      -\frac{4\rho^2}{k^3}g'(\theta_x)\zeta_x
      +\frac{\rho^2}{k^2}g''(\theta_x)\zeta_x^2
      +\frac{\rho^2}{k^2}g'(\theta_x)\eta_x\right]\dd x,
\end{aligned}
\]
with $\eta_0=0$. Since $\zeta_x=x+e_x$, the preceding estimate gives
\[
    \left\|\sup_{x\le L}|\zeta_x(k)|\right\|_{L^p(\Omega)}\le C_p.
\]
Using also that $h,h',h'',g,g',g''$ are bounded, BDG and H\"older's inequality imply, with
\[
    B_p(k):=\left\|\sup_{x\le L}|\eta_x(k)|\right\|_{L^p(\Omega)},
\]
that
\[
\begin{aligned}
B_p(k)
&\le C_p\left(k^{-3}+k^{-2}\left\|\sup_{x\le L}|\zeta_x|\right\|_{L^p(\Omega)}
        +k^{-1}\left\|\sup_{x\le L}|\zeta_x|^2\right\|_{L^p(\Omega)}\right) \\
&\quad +C_p k^{-1}
   \left\|\left(\int_0^L |\eta_s|^2\,\dd s\right)^{1/2}\right\|_{L^p(\Omega)} \\
&\quad +C_p\left(k^{-4}+k^{-3}\left\|\sup_{x\le L}|\zeta_x|\right\|_{L^p(\Omega)}
        +k^{-2}\left\|\sup_{x\le L}|\zeta_x|^2\right\|_{L^p(\Omega)}\right)
      +C_p k^{-2}\int_0^L\|\eta_s\|_{L^p(\Omega)}\,\dd s  \\
&\le C_p k^{-1}+C_p(k^{-1}+k^{-2})B_p(k).
\end{aligned}
\]
Here $p\ge2$ is used only to control $\|\sup_x|\zeta_x|^2\|_{L^p(\Omega)}$ by the fixed-frequency bounds in a higher moment. Enlarging $k_0$ if necessary, we absorb the last term and obtain
\begin{equation}\label{eq:eta-fixed-bound}
    \left\|\sup_{x\le L}|\partial_k^2\theta_x(k)|\right\|_{L^p(\Omega)}
    \le C_p k^{-1},
    \qquad k\ge k_0 .
\end{equation}

It remains to pass from deterministic $k$ to the supremum over $I_n$. We use the Banach-valued one-dimensional Sobolev inequality. If $f:I_n\to C[0,L]$ is absolutely continuous and $p>1$, then
\[
    \sup_{k\in I_n}\|f(k)\|_{C[0,L]}
    \le C
    \left(
      \left(\int_{I_n}\|f(k)\|_{C[0,L]}^p\,\dd k\right)^{1/p}
      +
      \left(\int_{I_n}\|f'(k)\|_{C[0,L]}^p\,\dd k\right)^{1/p}
    \right),
\]
where $C$ depends only on the length $|I_n|=2$. Apply this with
\[
    f(k)=\theta_\cdot(k)-k(\cdot),
    \qquad f'(k)=\partial_k\theta_\cdot(k)-\cdot.
\]
Taking $L^p(\Omega)$-norms and using Fubini's theorem, together with the fixed-frequency bounds, we obtain
\[
\begin{aligned}
&\left\|\sup_{k\in I_n}\sup_{x\le L}|\theta_x(k)-kx|\right\|_{L^p(\Omega)} \\
&\qquad \le C_p
    \left(\int_{I_n} k^{-p}\,\dd k\right)^{1/p}
    +C_p
    \left(\int_{I_n} k^{-p}\,\dd k\right)^{1/p}
    \le C_p n^{-1}.
\end{aligned}
\]
Applying the same Sobolev inequality with
\[
    f(k)=\partial_k\theta_\cdot(k)-\cdot,
    \qquad f'(k)=\partial_k^2\theta_\cdot(k),
\]
and using \eqref{eq:e-fixed-bound} and \eqref{eq:eta-fixed-bound}, we obtain
\[
    \left\|\sup_{k\in I_n}\sup_{x\le L}|\partial_k\theta_x(k)-x|\right\|_{L^p(\Omega)}
    \le C_p n^{-1}.
\]
This proves the lemma.
\end{proof}

\end{document}
